\begin{frame}{환경 고르는 실전 가이드 — 후보에 점수를 매기자}
  \small
  \begin{enumerate}
    \item \textbf{상태·행동 공간의 차원} — 차원이 높으면 학습 예산이
      급격히 커짐(예: HumanoidStandup은 $348\times17$).
    \item \textbf{보상 구조} — ``항상 음수(0이 목표)''형(Pendulum)과
      ``희소(목표 도달 시에만)''형(13.1 Reacher)은 \textbf{좋은 정책의
      모습} 자체가 달라 평가 지표의 의미가 달라진다.
    \item \textbf{시뮬레이션 격차가 보일 만큼 물리적으로 풍부한가} —
      13.1 도메인 무작위화를 제시하려면 Pendulum은 격차가 안 나와
      MuJoCo 로봇이 적합.
    \item \textbf{코드 재사용성} — Ch11.2의 PPO 코드가 거의 그대로
      적용되는가.
  \end{enumerate}
  \vskip0.5em
  \begin{block}{왜 점수 매기기를 하는가}
    기준을 먼저 정해두면 \textbf{팀 회의가 쉬워진다}(8.1절의 방식 그대로).
    환경은 ``임의로 고르는'' 것이 아니라 기준에 대해 점수를 받는 것.
  \end{block}
\end{frame}
